\section{Cross-paper consistency}\label{sec:audit}

This section adds no formal claim.  It checks that the way this paper
describes the companion papers agrees with what those papers state,
and that the two recognition sources of \Cref{sec:recognition} are
recorded as named premises rather than re-derived.

\subsection{Consistency of the framing}\label{sec:audit:cross-paper-consistency}

\paragraph{Recognition and derivation.}  In this paper a
recognition-mode landing and a derivation-mode landing are two
different kinds of conditional closure: the first names its
load-bearing source as a premise, the second derives its conclusion
from stated hypotheses.  Neither is treated as a weaker form of the
other.  This is consistent with the companion papers.  The
Riemann-hypothesis paper \cite{Tsiokos2026RH} proves the Riemann
hypothesis from a supplied Selberg recognition source and an
identification with the classical zeros; the P-versus-NP paper
\cite{Tsiokos2026PvNP} proves a fixed machine-level translation and
derives \(P\neq NP\) from an accepted negative source.  Both state these
premises as hypotheses.  The same separation is what keeps
\Cref{thm:attack-foreclosure-v5} from reading as a retraction of
\Cref{thm:attack-foreclosure-v4}: the first applies the second to
derivation routes only.

\paragraph{Column placement.}  The Riemann-hypothesis and P-versus-NP
substrates are placed in the trace-state-only column, and the
Navier--Stokes substrate in the witness-content-exposing column, by the
V-Differential classification of the hiddenness companion
\cite{Tsiokos2026CSL}.  That paper adopts the placements as
hypotheses, with an informal motivation for each; it does not prove
them.  This paper uses the same placements with the same status.  The
two trace-state-only substrates are mathematically unlike each other,
one analytic and one computational, but the criterion applied to them
is the same.

\paragraph{Source handling.}  In the P-versus-NP row the named source
is
recorded in the P-versus-NP paper \cite{Tsiokos2026PvNP}; in the
Riemann-hypothesis row the Selberg source is recorded in the
Riemann-hypothesis paper \cite{Tsiokos2026RH}.  In both rows the source is a named
premise and is not produced by a derivation in this paper.  A source
record is broader than the proposition it contributes: it carries its
provenance, the closure it belongs to and its column placement, while
the proposition is what enters the conditional argument.

\subsection{Recording sources}\label{sec:audit:source-record-discipline}

We use the following rule.  A recognition source must be named, its
substrate must be identified, and its role in a closure must be kept
separate from the proposition it contributes.  A proposed
recognition-mode landing that cannot name its source, identify its
column placement, or separate the source from the target proposition is
not counted as an instance of \Cref{thm:template-applicability}; such
cases belong to the open problems of \Cref{sec:predictions}.

For the P-versus-NP row, naming the source is bookkeeping, not
evidence for its negative content.  The hiddenness companion proves a
finite example in which two histories have different predictive classes
and the same current class while selected Boolean outputs read from the
two quotients agree; it does not connect its quotients with the classes
\(P\) and \(NP\).  The accepted negative content and the admission of
search events used for the standard complexity conclusion are separate
premises of the P-versus-NP paper.
